\section{A common solution of the tiling equations}\label{sec:model}

We now construct one graph satisfying all the equations of
\cref{sec:encoding,sec:activation}.  The useful output of each ordinary
channel will activate exactly the symbol $f(L_i(x))$ from
\cref{word:model}.  The high coordinates will record positions within
labelled blocks.  A further labelling of the horizontal residue set $S$
will let the two seed tests use a single common output $z$.

\subsection{Labelled blocks and their inverse map}

For every channel $i$ and digit $j\in J_i$, fix a partition of
$\Zmod{r_{ij}}$ into subsets of sizes $a_i$ and $b_i$, using the
nonnegative representation of $r_{ij}$ chosen in
\cref{sec:encoding}.  Label each subset of size $d\in\{a_i,b_i\}$
bijectively by $\Zmod d$.  Copy this same partition and these same labels
for every setting of the other digits in
$K_i=\prod_{j'\in J_i}\Zmod{r_{ij'}}$.  We obtain a partition
$\mathcal B_{i,j}$ of $K_i$ into labelled blocks.  A block fixes all
digits except $j$; its label records a position in its $j$-digit subset.
The blocks are subsets, not subgroups, and their labels need not respect
addition in $K_i$.

Define a map $C_{i,j}:K_i\to P_i$ and, for each integer $t$, a
permutation $T_{i,j}(t)$ of $K_i$ as follows.  If $k_i$ has label $y$ in
a block of size $d$, put
\[
 C_{i,j}(k_i)=
 \begin{cases}
 (y,0),&d=a_i,\\
 (0,y),&d=b_i.
 \end{cases}
\]
The permutation $T_{i,j}(t)$ sends this point to the point with label
$y+t$ in the same block, with addition in $\Zmod d$.
The map $C_{i,j}$ is independent
of every digit except $j$, since the partitions and labels were copied
unchanged over the other digits.  It is not independent of $j$: each
block has at least two positions with different useful outputs.

The following inverse calculation is the reason for the fibre
orderings in \cref{act:fibre-orderings}.

\begin{lemma}[Block inverse]\label{model:block-inverse}
Fix $i,j$, a target horizontal coordinate $x_1\in\Z$, and a target
low coordinate $k_i\in K_i$.  For each $(l,\delta)\in K_i\times\Delta_i$,
choose an integer $h_1(l,\delta)$ satisfying
\[
 h_1(l,\delta)\equiv l\pmod{r_i},\qquad
 h_1(l,\delta)\equiv\rho^a(\delta)\pmod{a_i},\qquad
 h_1(l,\delta)\equiv\rho^b(\delta)\pmod{b_i}.
\]
Then the map
\begin{equation}\label{model:block-map}
 (l,\delta)\longmapsto
 \left(
 C_{i,j}(k_i-l)+e_\delta,
 T_{i,j}\bigl(x_1-h_1(l,\delta)\bigr)(k_i-l)
 \right)
\end{equation}
is a bijection from $K_i\times\Delta_i$ to $P_i\times K_i$.
\end{lemma}

\begin{proof}
Fix a target $((\alpha^*,\xi^*),\beta_i^*)$.
The point $\beta_i^*$ specifies one block $B\in\mathcal B_{i,j}$ and
its label $q$.  As each $T_{i,j}(t)$ preserves its block, any preimage
must have its source $k_i'=k_i-l$ in $B$.

Suppose first that $B$ has size $a_i$.  Choose $\delta$ by the two
conditions
\begin{equation}\label{model:a-inverse}
 \begin{aligned}
 e_\delta^b&=\xi^* &&\text{in }\Zmod{b_i},\\
 e_\delta^a+\rho^a(\delta)&=\alpha^*+x_1-q
       &&\text{in }\Zmod{a_i}.
 \end{aligned}
\end{equation}
The first fibre ordering makes this choice unique.  Set
$y=\alpha^*-e_\delta^a$, let $k_i'$ be the unique point of $B$ with
label $y$, and set $l=k_i-k_i'$.  The useful output in
\eqref{model:block-map} is then $(\alpha^*,\xi^*)$.  The high output
has label
\[
 y+x_1-h_1(l,\delta)
 =\alpha^*+x_1-
   \bigl(e_\delta^a+\rho^a(\delta)\bigr)
 =q
 \quad\text{in }\Zmod{a_i},
\]
so it is $\beta_i^*$.  Conversely, both target coordinates force
exactly these choices of $\delta,y,k_i'$ and $l$.

If $B$ has size $b_i$, use the second fibre ordering to choose the
unique $\delta$ satisfying
\begin{equation}\label{model:b-inverse}
 \begin{aligned}
 e_\delta^a&=\alpha^* &&\text{in }\Zmod{a_i},\\
 e_\delta^b+\rho^b(\delta)&=\xi^*+x_1-q
       &&\text{in }\Zmod{b_i}.
 \end{aligned}
\end{equation}
Now $y=\xi^*-e_\delta^b$ specifies $k_i'$ in $B$, and again
$l=k_i-k_i'$.  The high label is
$\xi^*+x_1-(e_\delta^b+\rho^b(\delta))=q$ in $\Zmod{b_i}$.
These choices are necessary and sufficient, proving the bijection
in both cases.
\end{proof}

For an ordinary channel $i$, put $j_i(x)=f(L_i(x))$ and define
\begin{equation}\label{model:ordinary-outputs}
 c_i(x,k)=C_{i,j_i(x)}(k_i),\qquad
 \beta_i(x,k)=T_{i,j_i(x)}(x_1)k_i.
\end{equation}
Thus its active set is exactly $\{j_i(x)\}$, and $c_i$ has the
required dependence on $(L_i(x),k_i)$.

Fix a target base point $(x,k)$ in the activation test for this
channel.  Every batch has $h\in\Z u_i$, so its source $x'=x-h$
satisfies $L_i(x')=L_i(x)$.  Consequently all these sources use the
same block partition.  Their low coordinates are $k_i'=k_i-l$ and
their high outputs are
$T_{i,j_i(x)}(x_1-h_1)k_i'$.  The shear
\eqref{enc:shear} preserves this high output when the batch adds
$h_1$ to $v_i$.  Hence the batch values $(c_i,\beta_i)$ are precisely
\eqref{model:block-map}.  By \cref{model:block-inverse,act:fibre-criterion},
the ordinary activation equation is satisfied.  Its freely varied
$z$-coordinate places no condition on the common output to be defined
next.

\subsection{Seed regions and one common output}

Choose disjoint subsets $S_{\eta_1},S_{\eta_2}\subset S$ such that
\begin{equation}\label{model:seed-regions}
 \begin{gathered}
 \abs{S_{\eta_1}}=3,\qquad \abs{S_{\eta_2}}=4,\\
 S_{\eta_t}\subset
 \{(s_1,s_2)\in S:s_1\equiv0,\ s_2\equiv t\pmod p\}
 \qquad(t=1,2).
 \end{gathered}
\end{equation}
Each of the two specified residue classes has $p^2$ points, so these
choices are possible.  Write
$S_0=S\setminus(S_{\eta_1}\cup S_{\eta_2})$.  Since $p>200$ is prime,
$D=p^4\equiv1\pmod6$, and therefore $\abs{S_0}=D-7$ is divisible by
$6$.  Recall that $a_{\eta_1}=2$ and $a_{\eta_2}=3$.

Partition each of these three regions into label spaces according to
the following table.  A label space is simply a subset equipped with
the displayed bijection; these are not assertions about subgroups of
$S$.
\begin{center}
\begin{tabular}{llll}
\toprule
Region & Inactive seeds & Label space & Number of copies\\
\midrule
$S_{\eta_1}$ & $\eta_2$ & $\Zmod3$ & $1$\\
$S_{\eta_2}$ & $\eta_1$ & $\Zmod2$ & $2$\\
$S_0$ & $\eta_1,\eta_2$ & $\Zmod2\times\Zmod3$ & $(D-7)/6$\\
\bottomrule
\end{tabular}
\end{center}
For a seed $i$ inactive at $s$, denote its label coordinate by
$y_i(s)\in\Zmod{a_i}$.  Thus a point of $S$ is specified by its region,
the particular copy of the label space in that region, and its
inactive-seed labels.

For $t\in\Z$, define a permutation $R(t)$ of $S$ by preserving each
label space and adding $t$ to each of its label coordinates.  In
particular, $R(t)$ preserves all three regions, and its inverse is
$R(-t)$.  Fix an arbitrary bijection $\lambda:S\to\Zmod D$ and set
\begin{equation}\label{model:shared-output}
 z(x,k)=\lambda\bigl(R(x_1)s(x)\bigr).
\end{equation}
Neither $\lambda$ nor $R(t)$ is required to be a homomorphism.  The
set $S$ only indexes the values of this one cyclic output coordinate.

For a seed $i$, its digit set is the singleton $\{*\}$.  Abbreviate
$C_{i,*}$ and $T_{i,*}$ to $C_i$ and $T_i$.  Define
\begin{equation}\label{model:seed-outputs}
 (c_i(x,k),\beta_i(x,k))=
 \begin{cases}
 \bigl(C_i(k_i),T_i(x_1)k_i\bigr),&s(x)\in S_i,\\
 \bigl((y_i(s(x)),0),k_i\bigr),&s(x)\notin S_i.
 \end{cases}
\end{equation}
The useful output depends only on $(s(x),k_i)$, as required by
\eqref{enc:dependence}.  Its sole digit is active exactly when
$s(x)\in S_i$.

\subsection{The inverse for each seed test}

We verify both seed equations using the same functions
\eqref{model:shared-output} and \eqref{model:seed-outputs}.  Fix a seed
$i$, write $i'$ for the other seed, and fix a target base $(x,k)$.
The batch indexed by $(l,\delta,\tau)$ has a predetermined offset
$h=h(l,\delta,\tau)$.  Its source satisfies
\[
 x'=x-h,\qquad s'=s(x')=s(x)-\tau,\qquad k_i'=k_i-l.
\]
In addition to the congruences used by \cref{model:block-inverse},
the seed offsets satisfy
\begin{equation}\label{model:other-seed-congruence}
 h_1\equiv0\pmod{a_{i'}}.
\end{equation}
We will recover a unique batch from every target
\[
 \bigl(c_i^*,\beta_i^*,z^*\bigr)
 \in P_i\times K_i\times\Zmod D,
 \qquad c_i^*=(\alpha^*,\xi^*).
\]
By \cref{act:fibre-criterion}, this is exactly the required tiling
property.  Put $s''=\lambda^{-1}(z^*)$.  Because $R(x_1')$ preserves
each label space, the points $s'$ and $s''$ must be in the same
region and the same copy of its label space.

\paragraph{The seed is active at the source.}
Suppose $s''\in S_i$.  In this region the only label coordinate
belongs to the other seed $i'$.  Equation
\eqref{model:other-seed-congruence} gives
\[
 y_{i'}(s'')=y_{i'}(s')+x_1-h_1
            =y_{i'}(s')+x_1
 \quad\text{in }\Zmod{a_{i'}}.
\]
Thus $s'$ is uniquely determined by $s''$: preserve its label-space
identity and subtract $x_1$ from its label.  Equivalently,
$s'=R(-x_1)s''$ within this region.  We now know
$\tau=s(x)-s'$.  For this fixed $\tau$, apply
\cref{model:block-inverse} to the offsets $h_1(l,\delta,\tau)$ and the
fixed seed partition.  The appropriate formula
\eqref{model:a-inverse} or \eqref{model:b-inverse} recovers $\delta$;
the source label then fixes $k_i'$ and $l$.  The resulting source has the prescribed $s'$
and sends its $z$-coordinate to $z^*$.  Every target in this case
therefore has exactly one preimage.

\paragraph{The seed is inactive at the source.}
Suppose $s''\notin S_i$, and let $q=y_i(s'')$.  The source high
coordinate equals its low coordinate, so necessarily
\[
 k_i'=\beta_i^*,\qquad l=k_i-\beta_i^*.
\]
The target useful output and the rotated $i$-label impose
\[
 e_\delta^b=\xi^*,\qquad
 y_i(s')=\alpha^*-e_\delta^a,\qquad
 q=y_i(s')+x_1-\rho^a(\delta).
\]
Consequently $\delta$ is the unique index satisfying
\eqref{model:a-inverse}.  These equations then fix the source label
$y_i(s')$.  If the other seed is also inactive in this region, its
source label is uniquely fixed by
\[
 y_{i'}(s')=y_{i'}(s'')-x_1
 \quad\text{in }\Zmod{a_{i'}},
\]
using \eqref{model:other-seed-congruence}.  If the other seed is
active, there is no second label to recover.  Keeping the already
known label-space identity, these labels determine exactly one
$s'$, and then $\tau=s(x)-s'$ fixes the remaining batch index.
The offset for this recovered batch satisfies all the congruences
used in the calculation, so it produces exactly the desired target.
The recovery uses only the prescribed residues of $h_1$, so the offset indexed by the recovered $(l,\delta,\tau)$ has exactly the label action used to determine those indices. There is no dependence on a later choice of an integer lift. The necessity of every recovery step proves uniqueness.

For emphasis, on $S_0$ both labels are recovered on the same label
space.  In the test of $i$ their forward values are
\[
 \begin{aligned}
 y_i(s'')&=\alpha^*+x_1-
                  \bigl(e_\delta^a+\rho^a(\delta)\bigr),\\
 y_{i'}(s'')&=y_{i'}(s')+x_1.
 \end{aligned}
\]
The first line selects $\delta$ by the fibre ordering and then fixes
$y_i(s')$; the second fixes $y_{i'}(s')$ directly.  Interchanging $i$
and $i'$ proves the other seed equation with exactly the same
shared output $z$.  Each equation has its own unique representing
batch for a target; there is no requirement that those two batches
coincide.

\subsection{Compatibility of all the equations}

The definitions \eqref{model:ordinary-outputs},
\eqref{model:shared-output}, and \eqref{model:seed-outputs} specify
$c_i(x,k)$ and $\beta_i(x,k)$ for every channel, as well as $z(x,k)$,
at every base point $(x,k)$.  By
\eqref{enc:shear}, each chosen $\beta_i$ determines exactly one
$v_i\in\Zmod{r_i^2}$ with low coordinate $k_i$.  Taking these points
for all $(x,k)$ defines a single set $A\subset G$ satisfying the graph
equation $A\oplus H_0=G$.

The inverse calculations hold for every integer $x_1$, including
negative values, since all rotations use addition in the specified
finite label groups.  They use only the stated residues of $h_1$,
together with $s(h)=\tau$ in a seed test and $L_i(h)=0$ in an ordinary
test.  Thus any choices of the permitted Chinese remainder lifts,
including those made to separate distinct batches, give the same
verification.  The actual integer $x_1$ in the high outputs and in
$z$ causes no extra constraint: the dependence equations
\eqref{enc:dependence} restrict only the useful outputs $c_i$, not
$\beta_i$ or $z$.  The useful outputs have exactly those prescribed
dependencies, and all activation equations have just been checked
on this one graph.

It remains to check the exclusion equations.  At every $x$, the
unique active ordinary word is
\[
 \bigl(f(L_0(x)),\ldots,f(L_{N-1}(x))\bigr),
\]
which is allowed by \cref{word:model}.  Hence no forbidden ordinary
word has all its labels active.  If the seed $\eta_t$ is active at
$x$, then \eqref{model:seed-regions} gives
$L_n(x)\equiv t\pmod p$ for every ordinary $n$.  Since $t\in\{1,2\}$
is nonzero modulo $p$, we have $f(L_n(x))=t$.  No forbidden seed and
ordinary-symbol pair is simultaneously active either.
\Cref{enc:exclusion} now verifies every exclusion tile.

Combining this common solution with the obstruction already proved
in \cref{act:nonperiodic} yields the required intermediate result.

\begin{theorem}\label{model:system}
There are a finite cyclic group $V$ and finitely many nonempty finite
sets $F_1,\ldots,F_s\subset G=\Z^2\times V$ for which the simultaneous
tiling equations
\[
 A\oplus F_\nu=G\qquad(1\leq\nu\leq s)
\]
have a common solution, but no common solution is invariant under a
finite-index subgroup of $G$.
\end{theorem}

\begin{proof}
Use the graph, dependence, exclusion, and activation tiles
constructed in \cref{sec:encoding,sec:activation}.  They are finite
nonempty sets, and their finite coordinate group $V$ is cyclic.
The graph constructed in this section solves all their equations.
\Cref{act:nonperiodic} excludes every fully periodic common solution.
\end{proof}
